\section{Evaluation protocol}
\label{sec:protocol}

\subsection{Effective words per minute}
Error rates compare decoders but do not tell whether a user saves time. We therefore simulate a user who reads the system output, accepts or corrects a commit, picks from an offer, or types the sentence with a fallback keyboard. For item $i$ with intended sentence $y_i$, the simulation returns the total time $t_i$ (input, reading, selections, corrections and retyping) and the final text. The effective words per minute are
\begin{equation}
  \text{\ewpm{}} = 60 \cdot \frac{\sum_i \mathbf{1}[\text{final text}_i = y_i]\;\mathrm{words}(y_i)}{\sum_i t_i},
\end{equation}
a ratio of sums over all test items in which a sentence counts only if it is exactly right at the end. Like the BCI utility metric \citep{dalseno2010utility}, \ewpm{} charges the time needed to fix errors. A \emph{silent error} is a wrong final text that the user did not correct; it reaches the recipient and counts zero words. \ewpm{} does not price the harm of a silent error beyond the lost words, so we report the silent-error rate separately.

\subsection{User model and baselines}
The user model is a gaze-typing user. The fallback keyboard runs at 20 characters per minute, and scanning and selecting one of $k = 3$ word alternatives takes 4.5\,s; both are matched to measurements of one eye-gaze typing user with ALS \citep{cai2024gaze}. All other time costs are assumptions (\cref{tab:user}). The neural input runs at a device-specific speed: the measured 159.3 characters per minute of the intracortical participant; 75.8 characters per minute for sEMG (the 0.792\,s median press period); and 152 characters per minute for MEG (the healthy participants' typing speed).

\paragraph{Silent errors and the asymmetry between commit and offer.} After a \act{commit}, the user reads the sentence and misses each wrong word independently with probability 0.05; a missed wrong word becomes a silent error. After an \act{offer}, the user compares the candidates on purpose and sees every error, and chooses the cheapest of picking, fixing words or retyping, using the true costs. This asymmetry favours offers. It encodes the hypothesis that committed text is checked less carefully than offered text, which is not measured. The model does not make fluent wrong sentences harder to notice than garbled ones: a sentence with more wrong words is more likely to keep at least one of them. \Cref{sec:res-main} reports how the MEG conclusions change under other user-model settings. \todo{R1: user-model sensitivity grid on all three devices, a symmetric-offer variant and a misspecified-user evaluation without refitting.}

\paragraph{Population mismatch.} The fallback and selection costs are those of an eye-gaze user with ALS. This matches the intracortical participant, who has ALS. The MEG and sEMG decoders, however, read overt typing by able-bodied people, who would not need a gaze keyboard. For these two devices, \ewpm{} answers a hypothetical question: how fast would a gaze-typing user be with a decoder that has this error profile at this input rate?

\paragraph{Baselines.} \emph{Gaze typing only}: the same user types every sentence with the fallback keyboard and no BCI. This baseline is error-free and has no reading or correction time, which slightly favours it, but it also has no LM-based word prediction, which favours the BCI arms; LM prediction is known to speed up gaze typing \citep{ward2000dasher,cai2024gaze}. \todo{R1: LM-assisted gaze baseline (same 6-gram, word completion with $k = 3$ slots).} \emph{Always commit}: commit the top candidate of every item. \emph{Always offer}: offer the top candidates of every item and never commit; this is what most BCI spellers and AAC systems do.

\subsection{Datasets and splits}
\paragraph{Intracortical speech (Brain-to-Text '25).} The public copy-task file holds the online phoneme CTC logits, the cue sentences and the published LM-rescored sentences of participant T15, an implant recipient in the BrainGate2 clinical trial, for 1,718 trials over 15 days \citep{card2024accurate,card2025data}. We use the 1,628 trials with the 125,000-word vocabulary, which span 14 days; the 90 trials with a 50-word vocabulary are excluded. The file has no trial time stamps; we take the file order within each day as the trial order. We fit on the first half of each day (812 trials) and test on the second half (816 trials); one sentence occurs in both halves. Inside the fit half, the frame and sentence temperatures use the first half of the fit trials and the policy uses all of them. We evaluate two decoders: our own phoneme-to-word decoder (phoneme beam, CMUdict lexicon, Tatoeba word unigram and character LM; test WER 0.149) and the published sentence (test WER 0.043), which we gate as described in \cref{sec:trust}.

\paragraph{Forearm sEMG (Toronto).} The dataset holds isolated key presses from 19 participants over two days \citep{achituve2025emg,achituve2024data}. Because it has no continuous typing, we compose sentences from single presses, with oracle press timing and no inter-key context: one real press window per character, at the median press period. Sentences are held-out Tatoeba sentences that the LM did not see; exact overlap is excluded, but near-duplicates were not checked, and the style matches the LM corpus. Each participant contributes 30 fit sentences built from presses of the first day-2 recording of each letter, and up to 30 test sentences from later recordings (570 fit, 551 test). With about 10 presses per letter per pool, a press can appear in several test sentences, so sentence-level intervals are too narrow and participant intervals are primary. A CNN classifies each press. Three conditions vary how much of the user's data the decoder saw: \emph{session-calibrated} (day 1 plus half of the day-2 recordings, cross-fitted: the decoder that predicts the test presses is trained on the first recordings, and the fit sentences are predicted by the decoder trained on the later ones), \emph{cross-session} (day 1 only) and \emph{cross-subject} (4-fold leave-participants-out). Within the fit sentences, one half (A) fits the temperatures and the LM weight and the other half (B) fits the policy; CRC and LTT use the last 142 sentences of half B and the grid the first half.

\paragraph{MEG typing (SpanishBCBL).} Healthy participants typed memorised Spanish sentences in an MEG scanner \citep{levy2026brain2qwerty,spanishbcbl}. We use 19 participants and split by a hash of the sentence text into 101 training, 11 validation and 17 test sentences. The split is within participants: the same 19 people appear in training and test, with different sentences. No test sentence appears in training or in the LM text. The test set has 637 instances of the 17 sentences; all trust parameters are fitted on the 437 validation instances of 11 sentences. Our asynchronous ConvCTC decoder (convolutional front end and 4 Conformer layers of width 256 \citep{gulati2020conformer}) receives the MEG segment from 0.4\,s before the first key press to 0.5\,s after the last, its length and the participant index, but no key times at inference. It uses key times as auxiliary training labels, and the sentence bounds come from the keyboard log, so it is not a free-running decoder and cannot be trained for a user who cannot type (\cref{app:meg}). The trust harness decodes its streams with a fixed beam setting (LM weight 0.5, frame temperature fitted on validation) whose top-1 CER is 0.522, higher than the 0.475 of the beam setting tuned on validation CER (\cref{tab:decoders}). A synchronous decoder with oracle key times gives the same trust results (\cref{app:meg}).

\subsection{Controls}
\emph{Decoder controls} show that the decoders use the neural signal: shuffled MEG, zero input, and a decoder trained with permuted labels. \emph{Trust-layer controls} come in two classes. \emph{Signal-destroying controls} keep the stream format but remove the neural information: frame-shuffled streams, near-uniform streams, zero-input decoder output and a permuted-label decoder. A stream-only layer should not commit these. The \emph{intent-mismatch control} is the real stream of a different sentence by the same participant. It carries genuine neural evidence, only for a sentence the user did not intend now, so a stream-only layer cannot detect it.

\subsection{Statistics}
We report 95\% percentile bootstrap confidence intervals and paired differences on the same items. The resampling unit is stated with every interval: trials \texttt{[trial]} and days \texttt{\{day\}} for Brain-to-Text '25; sentences \texttt{[sent]} and participants \texttt{\{participant\}} for sEMG, where participants are primary; for MEG trust results, unique sentences (cluster bootstrap over the 17 test sentences), participants, and both at once (two-way), with the two-way interval primary for comparisons with gaze typing. MEG decoder CER intervals over instances are too narrow and are not used for inference. We name two primary endpoints per device: the wrong-commit rate of the trust layer and the paired \ewpm{} difference between the trust layer and gaze typing. All other intervals are descriptive and not corrected for multiplicity. Decoder settings, user-model values, the policy grid, $\alpha$, $\delta$ and the online step sizes were fixed in the experiment configurations before the evaluation runs. The ``noticed errors only'' online runs were added after the first run without changing any setting. Every decoder and trust run uses one seed and one fit/test split. \todo{R1: decoder seeds and repeated fit/test splits; sentence-cluster CIs for MEG CER.}
